\documentclass[twocolumn,11pt]{article}

\usepackage[margin=1in]{geometry}
\usepackage{amsmath,amssymb}
\usepackage{natbib}
\usepackage{graphicx}
\usepackage{hyperref}
\usepackage{booktabs}
\usepackage{xspace}

\newcommand{\dsa}{DSA-110\xspace}
\newcommand{\wsclean}{\textsc{WSClean}\xspace}
\newcommand{\casa}{\textsc{CASA}\xspace}
\newcommand{\everybeam}{\textsc{EveryBeam}\xspace}
\newcommand{\aoflagger}{\textsc{AOFlagger}\xspace}
\newcommand{\jybeam}{Jy\,beam$^{-1}$\xspace}

\title{A Continuum Imaging Pipeline for the DSA-110 Array}
\author{J.~T.~Faber}
\date{2026 March}

\begin{document}
\maketitle

\begin{abstract}
We describe the continuum imaging pipeline developed for the Deep Synoptic
Array 110-antenna station (\dsa) at the Owens Valley Radio Observatory.  The
pipeline converts raw correlator output into calibrated wide-field images,
constructs hourly epoch mosaics of the meridian drift strip, and performs
forced photometry against catalog positions to produce daily light curves
of $\sim\!10^{4}$ compact radio sources.  Because \dsa is a fixed meridian
transit instrument observing at $1.31$--$1.50$\,GHz, the pipeline must
accommodate a continuously drifting field of view, per-epoch gain
calibration against catalog sky models, and automated quality control
across $\sim\!100$ tiles per day.  We detail the design choices that follow
from these instrument constraints and present the overall data flow from
HDF5 visibilities to variability metrics.
\end{abstract}


% ----------------------------------------------------------------
\section{Introduction}\label{sec:intro}

The Deep Synoptic Array 110-antenna station \citep[\dsa;][]{hallinan2019}
is an L-band radio interferometer at the Owens Valley Radio Observatory
designed for the detection and localization of fast radio bursts and for
commensal continuum science.  The array comprises 117 antennas of 4.65\,m
diameter arranged in a core--outrigger configuration with baselines
from $\sim\!13$\,m to $\sim\!4.5$\,km.  As a fixed meridian transit instrument,
\dsa does not track sources; the sky instead drifts through the primary
beam at the sidereal rate.  This observing mode yields a natural daily
cadence well suited to monitoring flux variability on timescales of days
to months---the regime of extreme scattering events
\citep[ESEs;][]{fiedler1987}, refractive interstellar scintillation, and
intrinsic AGN variability.

The continuum pipeline described here was designed from the outset around
the constraints of a drift-scan array: each $\sim\!5$-minute transit
produces a single ``tile'' whose phase centre drifts in right ascension
over 24 correlator integrations, so the data require rephasing before
imaging; there is no dedicated calibrator scan, so gain calibration is
bootstrapped from catalog sky models; and the wide $4.1^{\circ}$ primary
beam demands direction-dependent beam correction.  The pipeline
automates the full path from raw HDF5 correlator dumps to per-epoch
forced-photometry light curves, with checkpoint recovery, flag-fraction
monitoring, and multi-catalog quality assurance at each stage.


% ----------------------------------------------------------------
\section{Instrument and Observing Mode}\label{sec:instrument}

Each correlator dump records 16 spectral windows (subbands) of 48
channels each, spanning 1311.25--1498.75\,MHz (187.5\,MHz total
bandwidth; channel width 244.14\,kHz).  The integration time is
12.885\,s per field, and each observation cycles through 24 fields
at successive meridian hour angles, producing a single tile of
duration $\sim\!310$\,s.  At a declination of $16\fdg1$ the 24 fields
span $\sim\!1\fdg5$ in right ascension, drifting at the sidereal rate
of $15\farcs03\,\mathrm{s}^{-1}\cos\delta$.

The primary beam full-width at half-maximum is approximately
$4\fdg1$ at 1.4\,GHz for the 4.65\,m dishes, closely matching the
image field of view adopted in the pipeline ($4800\times4800$ pixels
at $3''$/pixel, or $4\fdg0$ on a side).  The synthesized beam is
typically $60''\times34''$, set by the core--outrigger baseline
distribution.


% ----------------------------------------------------------------
\section{Pipeline Architecture}\label{sec:pipeline}

The pipeline proceeds in four stages: conversion, calibration, imaging,
and mosaicking with photometry (Figure~\ref{fig:flowchart} %conceptual
).  Each stage is
described below, with emphasis on the design choices motivated by the
transit observing mode.

\subsection{Conversion: HDF5 to Measurement Set}\label{sec:conversion}

Raw correlator output is stored as one HDF5 file per subband per
timestamp.  The 16 subbands for each observation are clustered within a
120\,s tolerance window and merged into a single Measurement Set (MS)
via \textsc{pyuvdata} \citep{hazelton2017}.

Three instrument-specific workarounds are required during conversion.
First, the \texttt{uvw\_array} in DSA-110 HDF5 files is natively
\texttt{float32}; it must be cast to \texttt{float64} before any
accumulation to avoid precision loss in the $u$, $v$, $w$ coordinates.
Second, the subbands arrive in \emph{descending} frequency order (sb00
$=$ highest), but multi-frequency synthesis imaging requires ascending
order; the pipeline sorts subbands and converts any negative channel
widths to positive (upper-sideband convention).  Third, the MS
\texttt{TELESCOPE\_NAME} must be set to \texttt{OVRO\_MMA} for \casa
compatibility but must be toggled to \texttt{DSA\_110} before any
\wsclean invocation so that \everybeam loads the correct Airy-disk beam
model.  Failure to toggle this field produces no runtime error but
introduces photometric errors of up to $\sim\!20\%$ near the primary
beam edge.

\subsection{Calibration}\label{sec:calibration}

Because \dsa has no slewing capability, there is no dedicated calibrator
scan.  Calibration is instead performed in two stages: a static bandpass
derived from a bright source transit (currently 3C\,454.3 at
$\sim\!12.5$\,\jybeam), and a per-epoch gain solution bootstrapped from
a catalog sky model.

\paragraph{Bandpass and static gain.}
The bandpass is solved per-SPW (i.e.\ not combining spectral windows)
because each of the 16 subbands exhibits a distinct instrumental
response shape.  Key parameters are a minimum signal-to-noise ratio of
5.0 per solution interval (elevated above the default of 3.0 to reject
noise-dominated solutions), a fill-gap interpolation width of 3 channels
($\sim\!730$\,kHz), and a minimum of 4 baselines per antenna.  An
initial phase-only pre-solve at 60\,s intervals removes
integration-timescale decorrelation before the bandpass fit.  Delay
($K$) calibration is omitted: empirical tests show residual phase slopes
$<3^{\circ}$ across the band, well within the tolerance for MFS imaging.

Static gain solutions are obtained in amplitude-plus-phase mode at two
timescales: a long solution interval (\texttt{solint=inf}, one solution
per $\sim\!5$-minute scan) for instrumental effects and a short interval
of 60\,s for atmospheric phase variations.  Application order is $B
\rightarrow G_{\mathrm{long}} \rightarrow G_{\mathrm{short}}$.  The
bandpass is valid for $\sim\!24$\,h, while gain solutions are valid for
$\sim\!1$\,h.

\paragraph{Per-epoch gain calibration.}
For science-quality imaging beyond the bandpass-calibrator transit, the
pipeline solves per-epoch gain tables against a unified catalog sky model
(FIRST, RACS, NVSS, and VLASS sources above 5\,mJy within $0\fdg3$ of
the field centre).  A preconditioner phase-only solve at 60\,s intervals
(combining SPWs to boost SNR by $\sim\!4\times$) provides an initial
estimate; a main solve then proceeds in two rounds: phase-only, then
amplitude-plus-phase.  If more than 30\% of solutions are flagged (an
indicator of insufficient SNR), the epoch falls back to the static daily
gain table with bandpass-only correction.

\subsection{Flagging}\label{sec:flagging}

The flagging strategy is a mandatory two-stage process.  Stage~1 applies
the \aoflagger \textsc{SumThreshold} algorithm via a custom Lua
strategy tuned for the OVRO radio-frequency interference (RFI)
environment: the base threshold is raised to 1.2 (from 1.0) with an
additional iteration for heavy RFI, and the transient threshold factor
is lowered to 0.8 for Starlink- and Iridium-band excision.  Stage~1
alone reduces the data to $\sim\!1\%$ flagged, but the residual kurtosis
remains $\sim\!536$ (far from Gaussian).  Stage~2 applies a $7\sigma$
median-absolute-deviation (MAD) clip on cross-correlation amplitudes,
bringing the total flagged fraction to $\sim\!2.4\%$ and the kurtosis
to $\sim\!3$.  Omitting Stage~2 leaves outlier rates at $\sim\!4\times$
the Gaussian expectation, which biases calibration solutions and
produces imaging artifacts.

\subsection{Imaging}\label{sec:imaging}

Each tile is phaseshifted to the median meridian position (the median RA
and Dec of the 24 fields, using a circular mean for RA to handle the
$0^{\circ}/360^{\circ}$ wrap) and imaged with \wsclean
\citep{offringa2014}.  The default imaging configuration uses Briggs
weighting with robust$\,=0.5$, an auto-masking threshold of $5\sigma$,
an auto-cleaning threshold of $1\sigma$, a major-cycle gain of 0.8, and
a UV minimum of $1\,k\lambda$ to exclude short-baseline emission.
Primary beam correction is applied post-imaging via \everybeam using an
Airy-disk model for the 4.65\,m dishes, with a primary beam limit of
0.2 (blanking pixels below 20\% of the peak response).

For survey-mode imaging targeting daily variability monitoring, the
configuration shifts to Briggs robust$\,=0.0$ (neutral weighting),
multiscale deconvolution with scales of $[0, 5, 15, 45]$ pixels to
recover both point and extended emission, and $n_{\mathrm{terms}}=2$
spectral-index fitting across the 18\% fractional bandwidth.

When the Image Domain Gridder \citep[IDG;][]{vandertol2018} is
selected, the 16 spectral windows must first be merged into a single SPW
(\casa \texttt{mstransform}).  This merge step resets the telescope name
to \texttt{OVRO\_MMA}; the pipeline automatically patches it back to
\texttt{DSA\_110} before each \wsclean invocation.

Optionally, a catalog-driven FITS mask (60$''$ circular apertures around
sources above a configurable flux threshold) is passed to \wsclean to
seed the clean, reducing the number of major cycles by a factor of
$2$--$4\times$.  When sky-model seeding is enabled, the catalog is
rendered to MODEL\_DATA via \wsclean \texttt{-predict} before the
imaging run.


% ----------------------------------------------------------------
\section{Mosaicking and Photometry}\label{sec:mosaic}

\subsection{Epoch Mosaics}\label{sec:epoch_mosaic}

Tiles are grouped into hourly epochs by UTC hour.  To ensure continuous
sky coverage and smooth flux transitions across epoch boundaries, each
epoch includes the last two tiles of the preceding hour and the first
two of the following hour, producing $\sim\!20$\,min of overlap.  Two
mosaicking tiers are available:

\begin{itemize}
  \item \textbf{Quicklook} (image-domain): pre-deconvolved tile FITS
  images are regridded to a common WCS via nearest-neighbor
  interpolation and coadded with weights
  $w = w_{\mathrm{rms}} \times w_{\mathrm{PB}}^{2}$,
  where $w_{\mathrm{PB}}$ is the Airy-disk primary beam response and
  the PB$^{2}$ factor implements variance weighting.
  Pixels below 10\% of the peak beam response are blanked.

  \item \textbf{Science/Deep} (visibility-domain): all tile MS files are
  phaseshifted to a common centre and jointly deconvolved with \wsclean
  using IDG and A-projection (\texttt{-grid-with-beam}) for
  direction-dependent beam correction.  This produces a properly
  deconvolved wide-field mosaic at the cost of significantly longer
  processing time.
\end{itemize}

\subsection{Forced Photometry and Variability}\label{sec:photometry}

Flux densities are measured at NVSS catalog positions using a
\citet{condon1997} matched-filter estimator: the PSF (a 2D Gaussian
derived from the \wsclean restoring beam) is convolved with the image,
and the weighted flux is
\begin{equation}
  S = \frac{\sum_{i} d_i\, K_i / \sigma_i^{2}}%
           {\sum_{i} K_i^{2} / \sigma_i^{2}},
\end{equation}
where $d_i$ is the image pixel value, $K_i$ the PSF kernel, and
$\sigma_i$ the local RMS from a background annulus or a BANE noise map.
Cross-matching uses a 30$''$ radius ($\lesssim 0.5$ synthesized beams)
and a minimum NVSS flux of 50\,mJy to ensure unconfused, high-SNR
comparisons.

Variability is quantified via the metrics of \citet{mooley2016}: the
reduced $\chi^{2}$ statistic $\eta$, the significance metric $V_{s}$,
and the modulation index $m$.  A source is flagged as a candidate
variable or ESE when its light curve exceeds a configurable
$\sigma$-threshold (default $5\sigma$, conservative).

\subsection{Quality Assurance}\label{sec:qa}

Automated QA operates at two levels.  Per-epoch, the pipeline checks
noise consistency (requiring $\mathrm{RMS}_{\max}/\mathrm{RMS}_{\min}
< 3$ across tiles) and gain-solution flag fractions ($<30\%$ for epoch
gains).  Per-day, the \textsc{verify\_sources} module cross-matches
DSA detections against NVSS and RACS catalogs and computes the median
DSA-to-catalog flux ratio; an epoch passes QA if this ratio exceeds 0.80
with at least three detections at $\mathrm{S/N}\geq3$.


% ----------------------------------------------------------------
\section{Automation and Fault Tolerance}\label{sec:automation}

The batch pipeline processes an entire day's observations
($\sim\!100$--200 tiles) through calibration, imaging, mosaicking, and
photometry in a single invocation.  Per-tile processing is parallelized
via Python's \texttt{ProcessPoolExecutor}, with a 30-minute timeout per
tile enforced by Unix signal alarms (necessitating process-level rather
than thread-level parallelism).  A JSON checkpoint file tracks completed
tiles so that a crashed run resumes from the last successful tile rather
than reprocessing the day.

The epoch gain calibration includes an automatic fallback: if the
catalog-bootstrapped gain solve yields more than 30\% flagged solutions,
the pipeline reverts to the static daily bandpass-only calibration for
that epoch, logging the fallback for downstream QA review.  Similarly,
if the input MS is more than 60\% flagged (e.g.\ due to heavy RFI),
self-calibration imaging is skipped and the catalog model is predicted
directly into MODEL\_DATA.

The end-to-end output is a set of hourly epoch mosaics, per-epoch
forced-photometry CSVs, and a summary QA table
(\texttt{qa\_summary.csv}) recording the median flux ratio, detection
count, gain-calibration mode, and pass/warn/fail status for each epoch.


% ----------------------------------------------------------------
\section{Summary}\label{sec:summary}

The \dsa continuum pipeline converts raw correlator output into
calibrated, mosaicked images and forced-photometry light curves with
daily cadence, targeting the detection of variable and transient compact
radio sources.  Its design is driven by the instrument's meridian transit
observing mode: catalog-bootstrapped gain calibration in lieu of
dedicated calibrator scans, per-epoch solutions to track hourly gain
drifts, automated rephasing of drifting field centres, and
direction-dependent primary beam correction for the wide $4^{\circ}$
field of view.  The pipeline has been validated on 3C\,454.3 commissioning
data (peak flux $12.5$\,\jybeam, consistent with catalog) and is
currently processing multi-day observing campaigns for ESE monitoring.


% ----------------------------------------------------------------
\bibliographystyle{aasjournal}
\begin{thebibliography}{9}

\bibitem[Condon(1997)]{condon1997}
Condon, J.~J. 1997, PASP, 109, 166

\bibitem[Fiedler et~al.(1987)]{fiedler1987}
Fiedler, R.~L., Dennison, B., Johnston, K.~J., \& Hewish, A. 1987,
Nature, 326, 675

\bibitem[Hallinan et~al.(2019)]{hallinan2019}
Hallinan, G., Ravi, V., Weinreb, S., et~al. 2019, BAAS, 51, 255

\bibitem[Hazelton et~al.(2017)]{hazelton2017}
Hazelton, B.~J., Jacobs, D.~C., Pober, J.~C., \& Beardsley, A.~P.
2017, JOSS, 2, 140

\bibitem[Mooley et~al.(2016)]{mooley2016}
Mooley, K.~P., Hallinan, G., Bourber, S., et~al. 2016, ApJ, 818, 105

\bibitem[Offringa et~al.(2014)]{offringa2014}
Offringa, A.~R., McKinley, B., Hurley-Walker, N., et~al. 2014,
MNRAS, 444, 606

\bibitem[van~der~Tol et~al.(2018)]{vandertol2018}
van~der~Tol, S., Veenboer, B., \& Offringa, A.~R. 2018, A\&A, 616,
A27

\end{thebibliography}

\end{document}
